\section{ViT: qué significa convertir una imagen en una secuencia}

Una vez establecidos la escala de datos y el transfer protocol, la clase llega a ViT. Su hipótesis de investigación no es que «attention siempre supera a convolution», sino esta: cuando hay suficientes datos y compute, ¿es posible reducir los priors visuales codificados de forma rígida, como la locality y la translation equivariance, para que el modelo aprenda la estructura espacial a partir de los datos? Por ello, la simplicidad de ViT es una ventaja y, al mismo tiempo, una exigencia de escala.

\subsection{Let the data speak: reducir la estructura visual diseñada a mano}

En una CNN, el local receptive field, el weight sharing y el submuestreo jerárquico forman un sesgo inductivo fuerte. ViT conserva menos estructura bidimensional: primero divide la imagen en bloques y los proyecta linealmente, y después procesa los token con un Transformer encoder genérico. La pregunta de investigación deja de ser «cómo diseñar un mejor módulo convolucional» y pasa a ser «si los datos pueden compensar los priors que se retiraron».

\lecturefigure{slide-16-letting-data-speak.jpg}{«ViT: letting the data speak» plantea la simplificación de la arquitectura como un experimento de escalamiento.}{Video oficial, 25:12.}

\begin{knowledgebox}{Lectura de la figura: menos priors no significa menos supuestos}
El patch size, la position embedding, el class token, la augmentation, el optimizer y la dataset composition siguen siendo decisiones de diseño de peso. Lo que se llama letting the data speak consiste en reducir, en términos relativos, la convolutional locality, no en eliminar el inductive bias. Unos priors más débiles suelen aumentar la flexibilidad de la representación, pero también amplían la sample complexity y la optimization burden.
\end{knowledgebox}

\lecturefigure{slide-17-cnn-vs-transformer.jpg}{La CNN codifica la estructura visual mediante una jerarquía de convoluciones; el Transformer ya demostró ser escalable en lenguaje, y ViT pregunta si puede aplicarse directamente a imágenes.}{Video oficial, 26:24.}

\begin{knowledgebox}{Lectura de la figura: lo que se compara son las rutas de información y la forma de compartir parámetros}
La convolución hace que los píxeles vecinos interactúen primero de manera local y amplía el receptive field a lo largo de las capas; la self-attention permite que cualquier patch establezca, en una sola capa, conexiones que dependen del contenido. La localidad de la CNN es más data-efficient; la interacción global de la attention es más flexible, pero su complejidad original crece de forma cuadrática con el número de token. No se trata de una dicotomía absoluta entre «local y global»: una CNN profunda también tiene un receptive field global, y ViT también aprende patterns locales.
\end{knowledgebox}

\subsection{Patch tokenization y class token}

Sea la imagen de entrada $x\in\mathbb{R}^{H\times W\times C}$ y sea $P$ el lado de cada patch; entonces el número de token es $N=HW/P^2$. Cada patch se aplana en un vector de dimensión $P^2C$ y se proyecta a $D$ dimensiones mediante la matriz lineal $E$. Después se antepone un class token aprendible y se suma la position embedding:
\[
z_0=[x_{\mathrm{cls}};x_p^1E;x_p^2E;\ldots;x_p^NE]+E_{\mathrm{pos}}.
\]
Aquí $x_p^i$ es el $i$-ésimo patch, $x_{\mathrm{cls}}$ sirve para agregar la información de clasificación y $E_{\mathrm{pos}}$ aporta el orden espacial de los token.

\lecturefigure{slide-18-vit-architecture.jpg}{ViT introduce la patch embedding, la position embedding y el class token en un Transformer encoder estándar.}{Video oficial, 28:32; Dosovitskiy et al. 2020.}

\begin{knowledgebox}{Lectura de la figura: el diagrama de la arquitectura se lee de abajo hacia arriba}
En la parte inferior, la imagen se divide en patches de tamaño fijo; la proyección lineal convierte cada patch en un token; la position embedding evita que la permutation invariance haga perder la información espacial; el encoder alterna multi-head self-attention y MLP; al final, el class token entra a la head de clasificación. El «Transformer estándar» de la figura sigue incluyendo LayerNorm, residual connection y el apilamiento de varias capas; no consiste en un único paso de attention.
\end{knowledgebox}

\begin{importantbox}{El costo computacional de la self-attention}
Para $N$ token y un hidden size $D$, el attention score $QK^{\top}$ requiere aproximadamente $O(N^2D)$ de cómputo y $O(N^2)$ de almacenamiento intermedio. Como $N=HW/P^2$, reducir a la mitad el lado del patch cuadruplica el número de token y multiplica por dieciséis el área de la attention matrix. Por eso, en tareas de visión de alta resolución no basta con fijarse en el número de parámetros: también hay que considerar el número de token y la activation memory.
\end{importantbox}

\begin{lstlisting}[caption={Patchify and prepend a class token}]
def vit_tokens(image, patch_size, projection, cls_token, pos_embed):
    patches = split_into_patches(image, patch_size)
    tokens = flatten(patches) @ projection
    tokens = concatenate([cls_token, tokens], axis=0)
    return tokens + pos_embed
\end{lstlisting}

\subsection{Resumen de la sección}

El experimento central de ViT consiste en debilitar los priors convolucionales, no en demostrar que una imagen «es, en esencia, lenguaje». El patchify determina el número de token y el compute, la position embedding restituye el orden y el class token ofrece una interfaz de agregación; una arquitectura sencilla traslada una mayor carga a los datos, la optimización y la escala.
